%% Generated by generate_slides.py. Edit this file or the source data, then run build.sh.
\begin{frame}{Why Obsidian?}
\subt{Obsidian edits an ordinary folder of Markdown files, rather than locking knowledge into an AI profile.}
\begin{columns}[c,onlytextwidth]
\column{.42\textwidth}\begin{itemize}\setlength{\itemsep}{.5em}\small
\item \textbf{Local files.} A vault is a directory; notes are human-readable .md files.
\item \textbf{Connections.} Links, backlinks, search and properties organize knowledge.
\item \textbf{Extensions.} Templates, Canvas and plugins add convenience.
\item \textbf{Important.} Obsidian is the editor; agents and protocols create the AI behavior.
\end{itemize}
\column{.55\textwidth}
\chgshoth{vault-interface.png}{3.35cm}{\chg{The lesson vault, a fake project on education and labour supply. Folder: lesson2/obsidian-vault.}}
\end{columns}
\end{frame}

\begin{frame}{Vaults, folders and plain Markdown}
\subt{\chg{Folders tell you where a file lives. A fixed set of files in every project (brief, state, decisions, next actions) tells the next session where the work stands.}}
\begin{itemize}\setlength{\itemsep}{0.85em}
\item \textbf{Vault.} Top-level folder, openable in Obsidian and another editor.
\item \textbf{Indexes.} Explicit maps tell readers and agents where project content lives.
\item \textbf{Metadata.} Markdown frontmatter can record type, owner, status and provenance.
\item \textbf{No requirement.} You do not need vector embeddings to create a useful first Brain.
\end{itemize}
\end{frame}

\begin{frame}{Links, backlinks and the Graph View}
\subt{Obsidian can reveal connections across notes and projects.}
\begin{center}\chgshoth[.86\textwidth]{vault-graph.png}{4.3cm}{}\end{center}
\tagline{\chg{The graph shows that two notes are linked. It cannot tell you whether the link makes sense or whether the note is still true: someone has to check.}}
\end{frame}

\begin{frame}{Canvas: think spatially}
\subt{A flexible surface for questions, evidence, methods and results.}
\begin{center}\chgshoth[.86\textwidth]{vault-canvas.png}{4.3cm}{}\end{center}
\tagline{\chg{Use Canvas to sketch. Once a choice is made, write it in DECISIONS.md, where the next session will look for it.}}
\end{frame}

\begin{frame}{Your first Brain: four useful project files}
\subt{Start with a small folder; let complexity emerge from recurring needs.}
\begin{columns}[c,onlytextwidth]\column{.565\textwidth}
\begin{itemize}\setlength{\itemsep}{.48em}\small
\item \textbf{PROJECT\_BRIEF.} What we are doing, why it matters, assumptions and scope.
\item \textbf{CURRENT\_STATE.} What is complete, what is blocked, which outputs are valid.
\item \textbf{NEXT\_ACTIONS.} Concrete, prioritized tasks that allow a cold restart.
\item \textbf{DECISIONS.} What is agreed, why, with dates and consequences.
\end{itemize}
\column{.405\textwidth}
\sboxblue{Add one index}{general/INDEX.md points to active projects and global rules.}
\sboxblue{Grow on demand}{SOURCES, AGENT\_MEMORY, LOG, wiki/ and skills/ only as real needs emerge.}
\end{columns}
\end{frame}

\begin{frame}{Connect an AI agent safely}
\subt{The agent needs explicit filesystem access and a repeatable entry point.}
\begin{columns}[c,onlytextwidth]\column{.565\textwidth}
\begin{itemize}\setlength{\itemsep}{.48em}\small
\item \textbf{Workspace.} Open the vault folder or provide the path using the agent’s supported tools.
\item \textbf{Entry point.} Give BOOTSTRAP.md: where to read, how to route, what not to touch.
\item \textbf{Permission scope.} Read-only for initial onboarding; writes require clear approval and Git diffs.
\item \textbf{Session protocol.} Start from index and state; end by updating validated durable memory.
\end{itemize}
\column{.405\textwidth}
\sboxblue{\chg{Terminal agents}}{\chg{Start Codex, Claude Code or OpenCode from inside the vault folder: \texttt{cd obsidian-vault} and then \texttt{codex}.}}
\sboxblue{\chg{Cursor}}{\chg{File, Open Folder, and pick the vault. The agent can now read every note in it.}}
\end{columns}
\tagline{\chg{Reading the vault is not remembering. The agent uses the Brain only if BOOTSTRAP.md tells it which files to open at the start and which to update at the end.}}
\end{frame}

\begin{frame}{Let AI build the first version}
\subt{\chg{Scaffolding a Brain is a small task, and the result is easy to test: open a new session and ask the agent where the project stands.}}
\begin{columns}[c,onlytextwidth]\column{.565\textwidth}
\begin{itemize}\setlength{\itemsep}{.48em}\small
\item \textbf{Provide constraints.} Work areas, privacy rules, naming, project file roles and desired workflows.
\item \textbf{Generate.} Index, template files, bootstrap, closeout and minimal validator.
\item \textbf{Review.} Inspect the folder tree, check descriptions and open the vault in Obsidian.
\item \textbf{Iterate.} Use an actual project for a week before adding complex memory backends.
\end{itemize}
\column{.405\textwidth}
\sboxblue{Simple prompt}{“Create a Markdown Brain with global index, project state, decisions and closeout; do not overwrite existing files.”}
\end{columns}
\medskip
\tagline{AI can create the structure; humans approve its memory policy.}
\end{frame}

\begin{frame}{Git + GitHub: version the knowledge}
\subt{Dropbox, Git and GitHub perform complementary jobs.}
\begin{columns}[c,onlytextwidth]
\column{.42\textwidth}\begin{itemize}\setlength{\itemsep}{.5em}\small
\item \textbf{Dropbox.} Synchronizes the same files across devices; not a substitute for reviewed commits.
\item \textbf{Git.} Records versions, shows diffs and allows local rollback.
\item \textbf{GitHub.} Optional remote host for collaboration and review; private by default for sensitive work.
\item \textbf{Security.} Exclude credentials; review changes before push; handle concurrent edits.
\end{itemize}
\column{.55\textwidth}
\chgshoth{github-commits.png}{4.9cm}{\chg{The commit history of this course on GitHub: every change has a date, a message and a diff.}}
\end{columns}
\end{frame}
